% Proof-ready supplement for Part B. The parent may renumber B.14a.
\begin{statement}{B.14a}{Proved}{An integral group lift of the exceptional-prime witness through class twelve}
Let $F=\langle x,y\rangle$ be the discrete free group,
$N_{12}=F/\gamma_{13}F$, and $z=(xy)^{-1}$.
There is an element $f\in\gamma_7N_{12}$ whose image in
$\gamma_7F/\gamma_8F$ is the specific integral vector
$q_7$ of B.12 and which satisfies the genuine group equations
\[
 f(y,x)f(x,y)=1,\qquad
 f(z,x)f(y,z)f(x,y)=1
 \quad\hbox{in }N_{12}.
\]
The same assertion holds with $q_7$ replaced by any integral
multiple. Consequently these finite nilpotent jets exist over
$\mathbb Z_2$ and $\mathbb Z_3$ without scaling $q_7$.
This is a finite-stage assertion; it does not assert an infinite
pro-$2$, pro-$3$, or full profinite lift.
\end{statement}
\begin{proof}
We use actual group commutators $[g,h]=ghg^{-1}h^{-1}$.
For a Lyndon word $w$, let $C_x=x$, $C_y=y$, and
$C_w=[C_u,C_v]$, where $w=uv$ is the standard Lyndon
factorization with $v$ the longest proper Lyndon suffix.
The leading Magnus term of $C_w-1$ is the standard integral
Lie bracket $b_w$. These leading terms form a basis of each
$\gamma_nF/\gamma_{n+1}F$. Induction over the lower central
series therefore gives unique integral coordinates for
$\gamma_7F/\gamma_{13}F$ in the ordered products of
$C_w$, $7\leq |w|\leq12$.

Set $x=1+X$, $y=1+Y$ in the truncated noncommutative
power-series ring over $\mathbb Z$. In particular
\[
 z=(1+Y)^{-1}(1+X)^{-1},
\]
and this full group series, rather than the linear expression
$-X-Y$ or logarithmic BCH coordinates, is used in every
substitution. If $A=g-1$, $B=h-1$, the exact identity
\[
 [g,h]-1=(AB-BA)(1+A)^{-1}(1+B)^{-1}
\]
constructs every needed Hall series by finite integer arithmetic.
The product of two terms of degree at least seven has degree at
least fourteen. Hence through degree twelve the Magnus expansion
of $\prod C_w^{a_w}$ is exactly
$1+\sum a_w(C_w-1)$, and both displayed group equations
are linear equations in all the Hall exponents.

We fix the seven weight-seven exponents to the coefficients of
$q_7$ and include \emph{every} exponent in weights eight through
twelve as an unknown. The simultaneous systems, retaining all
earlier free parameters, have the following exact ranks:
\[
\begin{array}{c|rrrrr}
\text{class}&8&9&10&11&12\\\hline
\text{unknown exponents}&30&86&185&371&706\\
\text{rank}&20&66&147&302&578\\
\text{free parameters}&10&20&38&69&128.
\end{array}
\]
These ranks and the final solutions are computed by
\texttt{code/e\_exceptional\_jets.py}; its full output and exact
coordinates are in \texttt{e\_exceptional\_jets\_12.json}.
At class twelve write the integral system as $Aa=b$.
Exact elimination at the primes $2,3,7$ gives solutions
$a_i=U_i/d_i$, with $U_i$ integral and
\[
 (d_2,d_3,d_7)
 =(43641883694233125,\ 720545000,\ 24).
\]
The corresponding denominator is a unit at the prime used in
each elimination. More strongly, the explicit identity
\[
 265181d_2-16061451206962d_3-1276d_7=1
\]
gives an integral solution
\[
 a=265181U_2-16061451206962U_3-1276U_7,
 \qquad Aa=b.
\]
The last stage's \texttt{combined\_jet.hall\_coordinates}
contains this vector. Together with the seven fixed coefficients
of $q_7$, it explicitly specifies the desired group element.

The program checks every associative-word coefficient of both
group defects through degree twelve, with exact rational or
integer arithmetic. Vanishing in the truncated Magnus ring
implies the asserted identities in $N_{12}$: if a group word
had a first nonzero lower-central component of degree at most
twelve, its nonzero Lyndon Lie polynomial would be its first
nonzero Magnus term. Conversely, a word in $\gamma_{13}F$
has no such term. This also justifies the independent full-word
verification of the recorded integral solution.

For completeness, selecting Lyndon-word rows when solving does
not discard an equation. Each cycle defect is primitive, up to
the given truncation, for the Magnus Hopf structure
$\Delta X=X\otimes1+1\otimes X+X\otimes X$.
Its first nonzero homogeneous part would be an ordinary free
Lie polynomial. The triangular Lyndon basis detects that part
on Lyndon words. Induction on degree proves that vanishing of
all selected homogeneous Lyndon coefficients is equivalent to
vanishing of the entire defect. In addition the code checks the
complete associative expansions of each final solution directly.

Finally, in this degree range replacing every Hall exponent by
its multiple $m$ multiplies both linear defects by $m$ and
replaces the leading term by $m q_7$. This proves the assertion
for all integral multiples, including the two multiples singled
out in the prompt. Passing to the pro-$p$ completion of $N_{12}$
or any finite $p$-group quotient preserves the verified equations.
\end{proof}

\paragraph{What remains of the exceptional-prime problem.}
The calculation evaluates the actual group-equation lifting
problem through the stated nilpotency class; it improves on the
mere existence of nonzero cohomological obstruction spaces and
on rational BCH jets. It does not prove that one of these jets
extends through every subsequent class. Existence at all classes,
with integral $p$-adic coordinates, would suffice: compactness of
the product of the Hall-coordinate copies of $\mathbb Z_p$
would give a compatible inverse-limit solution. The required
nonempty sets have been checked only to the finite bound above.
In particular the factor $6$ in an averaging argument is not a
unit for the two primes in question; taking norms may multiply
the first term at each new stage and therefore does not establish
an infinite lift of $q_7$, $2q_7$, or $3q_7$.
